%% =====================================================================
%% Cross-environment reproduction of a low-dose CT benchmark
%%
%% Target venue: Medical Physics (Technical Note)
%% Alternative:  Scientific Reports (if Nature Portfolio is preferred)
%% See ../PAPER_TARGET_JOURNALS.md row H.
%%
%% Every number in this manuscript comes from
%%   ../WS-1_dataset/R6_recalc/results/comparison_full764.json
%% via make_tables.py. Do not type numbers into the text; add them to the
%% generator. `python3 make_tables.py --check` fails when a table is stale.
%% =====================================================================
\documentclass[11pt]{article}
\usepackage[margin=1in]{geometry}
\usepackage{booktabs,amsmath,graphicx,hyperref,xcolor}
\usepackage[numbers,sort&compress]{natbib}
\newcommand{\todo}[1]{\textcolor{red}{[\textbf{TODO:} #1]}}

\title{\textbf{What agreement criterion does a reconstruction benchmark need?\\
A cross-environment reproduction of a low-dose CT evaluation}}

\author{\todo{author list, affiliations, ORCIDs, corresponding author}}
\date{\today}

\begin{document}
\maketitle

\begin{abstract}
\noindent\textbf{Purpose:} Benchmarks in low-dose CT reconstruction increasingly ship a numerical
agreement criterion so that a third party can verify a reported result. We asked what such a
criterion must look like to survive a reproduction in a different computing environment, and whether
a criterion that fails should be read as a failed reproduction.

\noindent\textbf{Methods:} An evaluation of five reconstruction methods over 764 held-out slices at
three dose levels was recomputed by an independent operator, in a separate environment differing in
operating system, CUDA build, and deep-learning framework version, from checkpoints and data
verified by hash. Agreement was assessed against the benchmark's original declared criterion, an
absolute difference of $10^{-6}$. Because that criterion failed, we tested the leading explanation
--- non-deterministic GPU kernels --- by re-running the affected methods with deterministic kernels
enabled, and separately audited the criterion itself across a ladder of absolute and relative
tolerances.

\noindent\textbf{Results:} The original declared criterion failed for three of five methods.
Re-running the same side with deterministic kernels reproduced the original bit-for-bit, and no
kernel non-determinism was observed in this environment. The criterion audit showed that none of
the tested absolute tolerances resolves the comparison --- it fails at $10^{-6}$ and equally at
$10^{-3}$ --- because the reported quantities span from a
dimensionless index near 0.9 to an observer statistic near $10^{6}$. A relative criterion at
$10^{-4}$, declared per metric in that metric's own units, admits every method, with a worst
observed relative difference of $6.36\times10^{-5}$.

\noindent\textbf{Conclusions:} A single absolute tolerance cannot express agreement across metrics of
differing scale, and within the range we tested loosening it does not repair it; only changing its
kind does. We recommend that
benchmarks declare agreement per metric in that metric's own units, and that a criterion failure be
investigated before it is either accepted as a discrepancy or relaxed.
\end{abstract}

\noindent\textbf{Keywords:} low-dose CT, reproducibility, benchmark, numerical agreement, detectability

\section{Introduction}

A benchmark that publishes a leaderboard implicitly promises that its numbers can be recomputed.
Increasingly that promise is made explicit: a released evaluation ships a tolerance, and a
recomputation that lands inside it is called a reproduction. The tolerance is rarely justified, and
it is usually a single number.

We report a reproduction in which that construction failed, and in which the failure was informative
rather than terminal. The original declared criterion of the benchmark under study was an absolute
difference of $10^{-6}$ applied uniformly. Three of five methods fell outside it. The obvious
reading --- that the recomputation disagreed with the release --- turned out to be wrong twice over:
no kernel non-determinism was observed as the cause in this environment, and none of the tested
absolute thresholds resolved the comparison.

The contribution of this note is not the verdict. It is the demonstration that a criterion's
\emph{kind} matters more than its magnitude, and a worked procedure for establishing that by
experiment rather than by preference.

\section{Methods}

\subsection{The evaluation under reproduction}

The benchmark evaluates five reconstruction methods --- RED-CNN~\citep{chen2017redcnn},
CTformer~\citep{ctformer2023}, LEARN~\citep{learn2018}, CoreDiff~\citep{corediff2023}, and a
Gaussian blur control included deliberately as a smoothing baseline --- over 764 held-out slices at
three simulated dose levels. Each run reports fidelity (PSNR, SSIM) and task-based detectability
(a channelized Hotelling observer AUC, a contrast-to-noise ratio, and a non-prewhitening observer
with an eye filter). Summary statistics and the number of contributing slices give 75 comparisons per method ---
5 seeds $\times$ 3 dose levels $\times$ 5 reported quantities --- and 375 in total across the
five methods.

The five quantities differ in scale by roughly six orders of magnitude, which is the crux of what
follows (Table~\ref{tab:scale}).

\begin{table}[ht]
\centering
\begin{tabular}{llll}
\toprule
Quantity & Units & Typical magnitude & Declared criterion \\
\midrule
PSNR & dB & $\sim$40--54 & relative, $10^{-4}$ \\
SSIM & index in [0,1] & $\sim$0.9 & relative, $10^{-4}$ \\
CNR & ratio & $\sim$2--7 & relative, $10^{-4}$ \\
CHO AUC & AUC in [0,1] & $\sim$1.0 & relative, $10^{-4}$ \\
NPWE & observer statistic ~1e5-1e6 & $1.6\times10^{5}$--$9.7\times10^{5}$ & relative, $10^{-4}$ \\
$n$ & count & 764 & absolute, exact \\
\bottomrule
\end{tabular}

\caption{\textbf{The reported quantities and their natural scales.} A single absolute tolerance is
simultaneously far too loose for a bounded index and far too tight for an observer statistic of
order $10^{5}$--$10^{6}$. The declared criterion in the right-hand column is the one adopted after
the audit in Section~\ref{sec:ladder}.}
\label{tab:scale}
\end{table}

\subsection{The two environments}

The reproduction was performed by an operator independent of the original run, using a separate
checkout, a separate virtual environment, and a separate output tree; no artifact of the original
run was overwritten. The recomputation environment was Python 3.12.10 with PyTorch 2.3.0+cu121 and
NumPy 1.26.4, on Windows with two NVIDIA RTX 4090 GPUs (driver 591.86). The reference environment,
in which the original results were produced, is the local run environment, and its exact version
record is the repository's preserved project-runtime record: Windows with two NVIDIA RTX 4090 GPUs
(driver 591.86), Python 3.12.10, PyTorch 2.3.0+cu121 (torchvision 0.18.0+cu121), NumPy 1.26.4,
SciPy 1.13.1, scikit-image 0.23.2, tifffile 2024.9.20 and mkl 2021.4.0. The complete 57-line
dependency freeze, including SHA-256-verified local wheel sources, is recorded in
\texttt{WS-1\_dataset/R6\_recalc/hashes/env\_pip\_freeze.txt}, and the source baseline is pinned at
commit \texttt{4875338}.

This is a reproduction across environments, not a re-seed within one. That distinction is the point:
a re-seed tests the pipeline's internal determinism, whereas a different library stack tests whether
the published number is a property of the method or of the machine.

\subsection{Provenance control}

Inputs were pinned by hash rather than by description: the source commit, with per-file SHA-256
verification of the six files carrying evaluation logic; the task specification; the five model
checkpoints; and every file of both data trees (364 and 25 entries). The recomputation environment
is recorded as a complete dependency freeze.

\subsection{The determinism ablation}
\label{sec:determinism}

The standard explanation for small cross-run differences on GPU is non-deterministic kernel
selection. This hypothesis is testable: with deterministic kernels
enabled, a re-run of the same side must be bit-identical, and any residual against the reference
observed under that configuration would not be attributable to kernel non-determinism in this
environment.

We re-ran the three methods outside the original declared criterion under the following settings, at a
cost of approximately 48 GPU-hours:

\begin{quote}
\begin{tabular}{@{}l@{}}
\texttt{cudnn.deterministic = True} \\
\texttt{cudnn.benchmark = False} \\
\texttt{torch.use\_deterministic\_algorithms(True)} \\
\texttt{CUBLAS\_WORKSPACE\_CONFIG = :4096:8} \\
\end{tabular}
\end{quote}

\subsection{The criterion audit}
\label{sec:ladder}

Independently of the verdict, we evaluated the comparison against a ladder of criteria: absolute and
relative, each at $10^{-6}$ through $10^{-3}$. The audit changes no measurement; it asks only which
criteria the \emph{same} recorded differences would satisfy.

\section{Results}

\subsection{Agreement}

Table~\ref{tab:agreement} reports the comparison. Under the original declared absolute criterion, 115
of 375 comparisons differed. The blur control and LEARN reproduced bit-for-bit; RED-CNN, CTformer
and CoreDiff did not.

\begin{table}[ht]
\centering
\begin{tabular}{lrrlr}
\toprule
Method & Comparisons & \multicolumn{1}{c}{$n_{\ne}$} & Worst relative difference & Outside \\
       &             & \multicolumn{1}{c}{(abs.\ $10^{-6}$)} & (metric)              & declared \\
\midrule
Blur (control) & 75 & 0 & bit-identical & 0 \\
LEARN & 75 & 0 & bit-identical & 0 \\
CTformer & 75 & 0 & bit-identical & 0 \\
RED-CNN & 75 & 45 & $6.36 \times 10^{-5}$, CNR & 0 \\
CoreDiff & 75 & 45 & $1.39 \times 10^{-5}$, CNR & 0 \\
\bottomrule
\end{tabular}

\caption{\textbf{Agreement by method.} $n_{\ne}$ counts comparisons differing at all under the
original declared absolute $10^{-6}$ criterion. The worst relative difference and its metric are given
for each method; the final column counts comparisons outside the declared per-metric relative
criterion of $10^{-4}$, which is zero everywhere. That two methods reproduce exactly, together with
the route-(a) rerun of Section~\ref{sec:determinism}, shows that the pipeline, data and code
snapshot are capable of bit-identical reproduction under this configuration; in this environment
the differences in the other three were not observed to depend on kernel non-determinism.}
\label{tab:agreement}
\end{table}

The largest single disagreement is instructive. For RED-CNN the worst absolute difference is 34.3,
which sounds disqualifying until one notes that it occurs in a quantity of magnitude
$\sim9\times10^{5}$ --- a relative difference of $6.36\times10^{-5}$.

\subsection{No kernel non-determinism was observed in this environment}

The prediction of Section~\ref{sec:determinism} was not supported in this environment. Re-running
the recomputation side with deterministic kernels reproduced the non-deterministic run exactly:
zero differing comparisons, bit-identical. No kernel non-determinism was observed in this
environment, so route (a) as scoped could not close the residual, and the difference against the
reference is treated as a systematic cross-environment offset.

This is the result that makes the rest of the note interpretable. Before it, environment and kernel
non-determinism were confounded, and any change to the criterion would have been an adjudication
made after seeing the outcome.

\subsection{None of the tested absolute tolerances resolves the comparison}

Table~\ref{tab:ladder} is the central result. The original declared criterion fails at $10^{-6}$; it
also fails at $10^{-5}$, $10^{-4}$ and $10^{-3}$, with the same three methods outside at every
level. Loosening an absolute tolerance by three orders of magnitude changes nothing within the
range tested, because the
binding comparison is on a quantity of order $10^{5}$--$10^{6}$, where an absolute budget of
$10^{-3}$ is still a relative demand of order $10^{-9}$.

A relative criterion behaves differently. It excludes two methods at $10^{-6}$ and $10^{-5}$ and
admits all five at $10^{-4}$.

\begin{table}[ht]
\centering
\begin{tabular}{llll}
\toprule
Criterion kind & Tolerance & Verdict & Methods outside \\
\midrule
Absolute & $10^{-6}$ & \textbf{FAIL} & CTformer, RED-CNN, CoreDiff \\
Absolute & $10^{-5}$ & \textbf{FAIL} & CTformer, RED-CNN, CoreDiff \\
Absolute & $10^{-4}$ & \textbf{FAIL} & CTformer, RED-CNN, CoreDiff \\
Absolute & $10^{-3}$ & \textbf{FAIL} & CTformer, RED-CNN, CoreDiff \\
\midrule
Relative & $10^{-6}$ & \textbf{FAIL} & RED-CNN, CoreDiff \\
Relative & $10^{-5}$ & \textbf{FAIL} & RED-CNN, CoreDiff \\
Relative & $10^{-4}$ & \textbf{PASS} & --- \\
Relative & $10^{-3}$ & \textbf{PASS} & --- \\
\bottomrule
\end{tabular}

\caption{\textbf{The criterion ladder.} Applied to the same recorded differences, with no
measurement changed. None of the absolute criteria tested is repaired by loosening it; only
changing its kind resolved the comparison among the criteria tested.}
\label{tab:ladder}
\end{table}

We note what the audit also establishes in the other direction: the original declared criterion was not
unattainable in principle. A 764-term float64 reduction accumulates a relative error of order
$10^{-13}$, so within a single environment the absolute $10^{-6}$ budget is comfortable. The
criterion was not too strict; it was the wrong kind.

\subsection{A provenance failure that numerical agreement would not have caught}

The first reproduction attempt disagreed with the reference by 9--13~dB in PSNR for CTformer. The
cause was not numerical: the operator guide named one checkpoint while the released run had used
another, a retrained variant. Re-running with the checkpoint the release actually used restored
agreement.

We report this because it is the failure mode a tolerance cannot detect. A criterion answers whether
two numbers agree; it is silent on whether they are numbers about the same thing. Hash-level
pinning of checkpoints, not their names, is what caught it.

\section{Discussion}

\subsection{Declare agreement per metric, in that metric's own units}

The recommendation follows from Table~\ref{tab:ladder} rather than from taste. A benchmark reporting
quantities of differing scale should declare, for each, whether agreement is absolute or relative
and at what tolerance, with counts required to match exactly. A single number applied across all of
them is a statement about the largest-magnitude quantity and is vacuous for the rest.

\subsection{A criterion failure is a finding, not a formality}

The sequence matters, and we would put it forward as a procedure. The criterion failed; the leading
explanation was stated as a testable prediction; the prediction was tested at real cost, and no
kernel non-determinism was observed in this environment; only then was the criterion amended
retrospectively, after the mismatch had been observed, and the superseded criterion and its original
verdict were retained verbatim beside the new one. The per-metric amendment is therefore a
retrospective adjustment based on the observed differences, not an independent or prospective
validation of the amended tolerance. Had the criterion been relaxed at the first failure, the
amendment would have been indistinguishable from adjusting a threshold until the result passed.

\subsection{Limitations}

Cross-environment agreement of a metric computation is not agreement of the underlying science.
What reproduces here is a computation, on fixed inputs, under declared hashes. Nothing in this note
speaks to whether the methods compared are clinically preferable, whether the detectability
surrogate used measures diagnostic content, or whether the dose levels are attainable in practice.

The comparison involves two environments. A third would establish whether the offset is systematic
across environments or particular to this pair, and this remains the clearest extension.
\todo{Decide whether to add a third environment; the constraint is moving checkpoints and data under
their use agreement, not compute.}

The dose levels are simulated rather than separately acquired, and the patient counts of the
underlying cohort are small; both limit what the reproduced numbers mean, though neither affects
whether they reproduce.

\section{Conclusions}

A reproduction that fails its original declared criterion has not necessarily failed. In the case
reported here the criterion was badly scaled, no kernel non-determinism was observed in this
environment, and the residual was a cross-environment offset four orders of magnitude inside a
defensible relative tolerance. Benchmarks that ship an agreement criterion should declare it per metric, in that
metric's own units, and should treat its failure as the beginning of an investigation.

\section*{Data and code availability}
\todo{Point at the released artifact, the comparison JSON, the re-derivation
script, and this note's table generator, once the repository's public release
decision is made.}

\section*{Author contributions}
\todo{CRediT statement.}

\section*{Competing interests}
\todo{Per-author declaration.}

\bibliographystyle{unsrtnat}
\bibliography{references}

\end{document}
